(1) Euler 方程在体固定主轴系是
\[
A\dot\omega_1=(A-C)\omega_2\omega_3,\quad
A\dot\omega_2=(C-A)\omega_3\omega_1,\quad
C\dot\omega_3=0.
\]
令 \(\Omega_b=(C-A)s/A\)，由初值 \(\omega_3=s\) 和前两式得到
\[
\omega_1(t)=w\cos(\Omega_bt),\qquad
\omega_2(t)=w\sin(\Omega_bt),\qquad
\omega_3(t)=s.
\]
\(\Omega_b\) 可以有符号，表示角速度横向分量在体坐标中的转动方向，不能将其绝对值误认为空间进动速度。

(2) 刚体动能
\[
E=\frac12A(\omega_1^2+\omega_2^2)+\frac12C\omega_3^2
=\frac12(Aw^2+Cs^2)
\]
恒定；体坐标角动量为 \((A\omega_1,A\omega_2,Cs)\)，故
\[
L=\sqrt{A^2w^2+C^2s^2}
\]
也恒定。因为空间角动量与对称轴的点积是 \(Cs\)，夹角 \(\vartheta\) 满足 \(\cos\vartheta=Cs/L\)。该夹角与时间无关。

(3) 在空间向量形式中，由对称惯量张量
\[
\boldsymbol L=A\boldsymbol\omega+(C-A)s\boldsymbol e_3,
\quad\Rightarrow\quad
\boldsymbol\omega=\frac{\boldsymbol L-(C-A)s\boldsymbol e_3}{A}.
\]
刚体轴的运动学为 \(\dot{\boldsymbol e}_3=\boldsymbol\omega\times\boldsymbol e_3\)，故
\[
\dot{\boldsymbol e}_3=\frac{\boldsymbol L}{A}\times\boldsymbol e_3 .
\]
无外力矩下 \(\boldsymbol L\) 在惯性系是固定向量，因而 \(\boldsymbol e_3\) 绕其方向以角速率
\[
\Omega_{\rm space}=L/A
\]
匀速沿圆锥面进动；角锥半角就是 \(\vartheta\)。

(4) 当 \(C\to A\)，体坐标频率 \(\Omega_b\to0\)，角速度在体内固定，符合球形刚体的自由自旋；若 \(w>0\)，所选材料对称轴仍可绕固定角动量方向旋转，描述一个运动学圆锥而不是动力学力矩进动。若 \(w\to0\)，\(\vartheta\to0\)，对称轴恰与角动量方向重合且在空间固定；此时尽管公式 \(L/A\) 有有限值，圆锥退化为一个点，所谓进动的方位角不再可测量。
